
% Copyright (c) 2015 - 2021 Mario Mlačak, mmlacak@gmail.com
% Licensed and published as Public Domain work.

% One chapter =========================================================
\chapter*{One}
\addcontentsline{toc}{chapter}{One}
\label{ch:One}

\begin{flushright}
\parbox{0.8\textwidth}{
\emph{God is not external to anyone, but is present with all things, though
they are ignorant that he is so. \\
\hspace*{\fill}{\textperiodcentered \textperiodcentered \textperiodcentered \hspace*{0.2em} Plotinus} } }
\end{flushright}

\noindent
One is chess variant which is played on 26 x 26 board, with white and
darker violet fields, and with light purple and fuchsia pieces. Star
colors are reversed colors of ordinary pieces, i.e. fuchsia and light
purple. In algebraic notation, columns are enumerated from 'a' to 'z',
and rows are enumerated from '1' to '26'. A new piece is introduced,
Starchild.

\clearpage % ..........................................................
% Starchild ***********************************************************

\section*{Starchild}
\addcontentsline{toc}{section}{Starchild}
\label{sec:One/Starchild}

% \vspace*{-1.5\baselineskip}
\noindent
\begin{wrapfigure}[11]{l}{0.4\textwidth}
\centering
\includegraphics[width=0.4\textwidth, keepaspectratio=true]{pieces/18_starchild.png}
\caption{Starchild}
\label{fig:18_starchild}
\end{wrapfigure}
Starchild cannot capture any piece, cannot check or checkmate opponent's King.
Starchild is celestial piece, it can participate in demoting-to-Pawn syzygy.
Starchild can be demoted to Pawn.

Starchild cannot be converted, but can be activated. If activated, it does not spend
momentum while moving. Starchild can activate own Wave and Starchild on step-fields.
Starchild can activate any own piece (except King), opponent's Starchild and any Star
on its miracle-fields.

% \vspace*{0.15\textheight}
\noindent
\begin{wrapfigure}[11]{l}{0.4\textwidth}
\centering
\includegraphics[width=0.4\textwidth, keepaspectratio=true]{pieces/star/22_one.png}
\caption{Star}
\label{fig:star/22_one}
\end{wrapfigure}
Starchild can't teleport. Starchild moves from starting to destination field in
opposite color in one step, without interacting with any piece on chessboard.

Starchild can resurrect any captured piece, except Kings, Stars, Monoliths. Waves
and Starchilds can be resurrected without resurrecting Starchild being oblationed.
Starchild can take any piece, except Kings, Waves, Stars, Starchilds and Monoliths,
for a non-interactive, viewing-only trance-journey.

Star colors in this variant are opposite colors of normal pieces.
In algebraic notation, symbol for Starchild is 'I'.

\clearpage % ..........................................................
% Movement ------------------------------------------------------------

\subsection*{Movement}
\addcontentsline{toc}{subsection}{Movement}
\label{sec:One/Starchild/Movement}

\vspace*{-1.1\baselineskip}
\noindent
\begin{figure}[!h]
\includegraphics[width=1.0\textwidth, keepaspectratio=true]{examples/22_o/scn_o_01_starchild_movement.png}
\caption{Starchild movement}
\label{fig:scn_o_01_starchild_movement}
\end{figure}

Starchild can move to any empty field in opposite color to starting location in
one step. Starchild is not hampered by any piece between starting and destination
field.

Here, light Starchild in the middle moved from its starting position. It is now
positioned at dark field, so in next move it can access any empty light field, in
a single step.

\clearpage % ..........................................................

\subsubsection*{Activating on step-fields}
\addcontentsline{toc}{subsubsection}{Activating on step-fields}
\label{sec:One/Starchild/Movement/Activating on step-fields}

\vspace*{-1.1\baselineskip}
\noindent
\begin{figure}[!h]
\includegraphics[width=1.0\textwidth, keepaspectratio=true]{examples/22_o/scn_o_02_starchild_activating_own_piece_init.png}
\caption{Activating Wave}
\label{fig:scn_o_02_starchild_activating_own_piece_init}
\end{figure}

Starchild can only activate own Waves and Starchilds on its step-fields, with 1
momentum. Here, both light Waves and own Starchild can be activated. Neither light
Pyramid nor any of other opponent's pieces can be activated; some of them are also
on the same color field as activating Starchild.

\clearpage % ..........................................................

\vspace*{-2.1\baselineskip}
\noindent
\begin{figure}[!h]
\includegraphics[width=1.0\textwidth, keepaspectratio=true]{examples/22_o/scn_o_03_starchild_activating_own_piece_end.png}
\vspace*{-1.3\baselineskip}
\caption{Wave activated}
\label{fig:scn_o_03_starchild_activating_own_piece_end}
\end{figure}

\vspace*{-0.3\baselineskip}
Activated Wave moves the same as Starchild, i.e. to any field in opposite color to
its starting position. There it can activate only own Starchilds, own and opponent's
Waves, with 1 received momentum. \\
Here, light Wave A is now activated, and it can only activate dark Wave. It cannot
activate opponent's pieces (dark Knight, Pegasus). Own pieces (light Starchild and
Wave B) can't be activated because their field is in the same color as Wave's
activation field.

\clearpage % ..........................................................

\subsubsection*{Activating Starchild}
\addcontentsline{toc}{subsubsection}{Activating Starchild}
\label{sec:One/Starchild/Movement/Activating Starchild}

\vspace*{-1.3\baselineskip}
\noindent
\begin{figure}[!h]
\includegraphics[width=1.0\textwidth, keepaspectratio=true]{examples/22_o/scn_o_04_activating_starchild.png}
\caption{Activating Starchild}
\label{fig:scn_o_04_activating_starchild}
\end{figure}

\hyperref[fig:10_wave]{Like Wave}, activated Starchild does not spend received
momentum for moving; if Starchild activates a piece, it transfers all of received
momentum to it. Wave activated by Starchild on its step-field moves (and acts) as
one; so, it can activate only own Starchilds, own and opponent's Waves on its
step-fields, and cannot activate any other pieces, neither own nor opponent's.

\clearpage % ..........................................................

\subsubsection*{Miracle-fields}
\addcontentsline{toc}{subsubsection}{Miracle-fields}
\label{sec:One/Starchild/Movement/Miracle-fields}

% \vspace*{-0.9\baselineskip}
\noindent
\begin{wrapfigure}[5]{l}{0.4\textwidth}
\centering
\includegraphics[width=0.192307692\textwidth, keepaspectratio=true]{examples/22_o/scn_o_05_miracle_fields.png}
\caption{Miracle-fields}
\label{fig:scn_o_05_miracle_fields}
\end{wrapfigure}
Miracle-fields are all fields immediately surrounding a piece horizontally,
vertically and diagonally. They are the same as step-fields of a King.

\vspace*{2.1\baselineskip}
\subsubsection*{Activating on miracle-fields}
\addcontentsline{toc}{subsubsection}{Activating on miracle-fields}
\label{sec:One/Starchild/Movement/Activating on miracle-fields}

% \vspace*{-0.9\baselineskip}
\noindent
\begin{wrapfigure}[8]{l}{0.4\textwidth}
\centering
\includegraphics[width=0.192307692\textwidth, keepaspectratio=true]{examples/22_o/scn_o_06_starchild_activating_on_miracle_fields.png}
\caption{Activating piece}
\label{fig:scn_o_06_starchild_activating_on_miracle_fields}
\end{wrapfigure}
Fields at which Starchild can activate a piece are miracle-fields; pieces that
can be activated are all own pieces (except King), and opponent's Starchild.

Note, Starchild cannot move to empty miracle-fields, if they're not in opposite
color to starting field; it can only activate a piece on its miracle-field.

Here, Starchild's activation fields are enumerated. Opponent's Bishop and own King
can't be activated, so only own Pyramid can be, with 1 momentum.

\clearpage % ..........................................................

% \vspace*{1.1\baselineskip}
\subsubsection*{Starchild cannot teleport}
\addcontentsline{toc}{subsubsection}{Starchild cannot teleport}
\label{sec:One/Starchild/Movement/Starchild cannot teleport}

% \vspace*{-0.9\baselineskip}
\noindent
\begin{wrapfigure}[6]{l}{0.4\textwidth}
\centering
\includegraphics[width=0.346153846\textwidth, keepaspectratio=true]{examples/22_o/scn_o_07_starchild_not_moving_monolith_init.png}
\caption{Moving into a Monolith}
\label{fig:scn_o_07_starchild_not_moving_monolith_init}
\end{wrapfigure}
Starchild cannot teleport. Starchild trying to capture either step- or miracle-field
at which Monolith (or a Star) is located, will emerge on an empty portal-field
surrounding that same Monolith (or a Star).

\vspace*{5.1\baselineskip}
\noindent
\begin{wrapfigure}[8]{l}{0.4\textwidth}
\centering
\includegraphics[width=0.346153846\textwidth, keepaspectratio=true]{examples/22_o/scn_o_08_starchild_not_moving_monolith_end.png}
\caption{Moving out of a Monolith}
\label{fig:scn_o_08_starchild_not_moving_monolith_end}
\end{wrapfigure}
In previous example Starchild moved into a Monolith; here, it emerges on an empty
portal-field of that same Monolith.

If there were no empty portal-fields, Starchild would be oblationed, i.e. removed
from chessboard as if captured by opponent.

\clearpage % ..........................................................

% \vspace*{2.1\baselineskip}
\subsubsection*{Moving a Star}
\addcontentsline{toc}{subsubsection}{Moving a Star}
\label{sec:One/Starchild/Movement/Moving a Star}

% \vspace*{-0.9\baselineskip}
\noindent
\begin{wrapfigure}[5]{l}{0.4\textwidth}
\centering
\includegraphics[width=0.192307692\textwidth, keepaspectratio=true]{examples/22_o/scn_o_09_starchild_moving_star_init.png}
\caption{Moving into a Star}
\label{fig:scn_o_09_starchild_moving_star_init}
\end{wrapfigure}
Starchild can activate a Star the same way as any other piece, i.e. by capturing
miracle-field at which Star is located. Activated Star receives 1 momentum.

\vspace*{2.1\baselineskip}
\noindent
\begin{wrapfigure}[4]{l}{0.4\textwidth}
\centering
\includegraphics[width=0.192307692\textwidth, keepaspectratio=true]{examples/22_o/scn_o_10_starchild_moving_star_end.png}
\caption{Star moving}
\label{fig:scn_o_10_starchild_moving_star_end}
\end{wrapfigure}
Once activated, Star can move to any empty miracle-field of activating Starchild,
which all are enumerated in example on the left.

\vspace*{2.3\baselineskip}
\noindent
\begin{wrapfigure}[8]{l}{0.4\textwidth}
\centering
\includegraphics[width=0.192307692\textwidth, keepaspectratio=true]{examples/22_o/scn_o_11_starchild_moving_star_activating.png}
\caption{Activating Starchild}
\label{fig:scn_o_11_starchild_moving_star_activating}
\end{wrapfigure}
Note, even if activated Starchild received more than 1 momentum, Star can move for
only one step.

Here, Star received all of initial 3 momentum gathered by the Rook, since neither
Wave nor Starchild expend momentum for its movement. Nevertheless, activated Star
can move for only one field.

Note that Star movement is optional, Starchild could emerge on an empty portal-field
around Star in question; if there are none, Starchild would be oblationed.

\clearpage % ..........................................................

% \vspace*{2.1\baselineskip}
\subsubsection*{Star movement blocked}
\addcontentsline{toc}{subsubsection}{Star movement blocked}
\label{sec:One/Starchild/Movement/Star movement blocked}

% \vspace*{-0.9\baselineskip}
\noindent
\begin{wrapfigure}[4]{l}{0.4\textwidth}
\centering
\includegraphics[width=0.192307692\textwidth, keepaspectratio=true]{examples/22_o/scn_o_12_star_movement_blocked_init.png}
\caption{Activating Star}
\label{fig:scn_o_12_star_movement_blocked_init}
\end{wrapfigure}
If there is no empty field around Star at actual activation, Star is not activated
(it remains at the same field), and Starchild is oblationed.

\vspace*{4.1\baselineskip}
\noindent
\begin{wrapfigure}[7]{l}{0.4\textwidth}
\centering
\includegraphics[width=0.192307692\textwidth, keepaspectratio=true]{examples/22_o/scn_o_13_star_movement_blocked_end.png}
\caption{Star blocked}
\label{fig:scn_o_13_star_movement_blocked_end}
\end{wrapfigure}
Here, Starchild (now "it-the-air") is about to activate Star; grey arrows show path
travelled over by the piece they point to. There is no empty field around Star, so it
cannot be activated; instead, Starchild is oblationed.

\clearpage % ..........................................................

% \vspace*{1.1\baselineskip}
\subsubsection*{Conversion immunity}
\addcontentsline{toc}{subsubsection}{Conversion immunity}
\label{sec:One/Starchild/Movement/Conversion immunity}

\vspace*{-1.4\baselineskip}
\noindent
\begin{figure}[!h]
\includegraphics[width=1.0\textwidth, keepaspectratio=true]{examples/22_o/scn_o_14_starchild_conversion_immunity_init.png}
\caption{Conversion immunity}
\label{fig:scn_o_14_starchild_conversion_immunity_init}
\end{figure}

\hyperref[sec:Mayan Ascendancy/Pyramid/Conversion]{Conversion} is a move in which
activated Pyramid reaches opponent's piece, if it's not King, on own side of board.
Pyramid is then oblationed, and reached piece is replaced by the same piece in own
color. Starchild cannot be converted, instead, original Starchild remains on
chessboard; conversioning Pyramid is still oblationed.

% ------------------------------------------------------------ Movement
\clearpage % ..........................................................
% Activating Wave -----------------------------------------------------

% \vspace*{-0.1\baselineskip}
\subsection*{Activating Wave}
\addcontentsline{toc}{subsection}{Activating Wave}
\label{sec:One/Starchild/Activating Wave}

% \vspace*{-0.9\baselineskip}
\noindent
\begin{wrapfigure}[8]{l}{0.4\textwidth}
\centering
\includegraphics[width=0.192307692\textwidth, keepaspectratio=true]{examples/22_o/scn_o_15_starchild_activating_wave_on_miracle_fields.png}
\caption{Activating Wave}
\label{fig:scn_o_15_starchild_activating_wave_on_miracle_fields}
\end{wrapfigure}
Wave activated by Starchild on its miracle-fields can activate a piece, with
1 momentum; any own piece (except King), and opponent's Waves can be activated.

Wave can also move to any empty miracle-field, regardless of color.

\subsubsection*{Wave cannot move a Star}
\addcontentsline{toc}{subsubsection}{Wave cannot move a Star}
\label{sec:One/Starchild/Activating Wave/Wave cannot move a Star}

\noindent
\begin{wrapfigure}[9]{l}{0.4\textwidth}
\centering
\includegraphics[width=0.192307692\textwidth, keepaspectratio=true]{examples/22_o/scn_o_16_starchild_activated_wave_not_moving_a_star.png}
\caption{Not moving a Star}
\label{fig:scn_o_16_starchild_activated_wave_not_moving_a_star}
\end{wrapfigure}
Wave cannot move a Star, even if activated by Starchild on its miracle-field.
Instead, Wave would emerge on any empty portal-field around a Star which it tried
to activate (here, fields 1 and 2). If there are no empty portal-fields, Wave is
oblationed.

\clearpage % ..........................................................

\subsubsection*{Wave cannot teleport}
\addcontentsline{toc}{subsubsection}{Wave cannot teleport}
\label{sec:One/Starchild/Activating Wave/Wave cannot teleport}

\vspace*{-0.1\baselineskip}
\noindent
\begin{wrapfigure}[4]{l}{0.4\textwidth}
\centering
\includegraphics[width=0.307692308\textwidth, keepaspectratio=true]{examples/22_o/scn_o_17_starchild_activated_wave_not_teleporting_init.png}
\vspace*{-0.4\baselineskip}
\caption{Moving into a Star}
\label{fig:scn_o_17_starchild_activated_wave_not_teleporting_init}
\end{wrapfigure}
Wave activated by Starchild cannot teleport, regardless if Wave was activated on
Starchild's step- or miracle-field.

\vspace*{5.7\baselineskip}
\noindent
\begin{wrapfigure}[7]{l}{0.4\textwidth}
\centering
\includegraphics[width=0.307692308\textwidth, keepaspectratio=true]{examples/22_o/scn_o_18_starchild_activated_wave_not_teleporting_end.png}
\vspace*{-0.4\baselineskip}
\caption{Moving out of a Star}
\label{fig:scn_o_18_starchild_activated_wave_not_teleporting_end}
\end{wrapfigure}
Instead of teleporting, Wave emerges on an empty portal-field around Monolith or a
Star through which it tried to teleport.

If there is no empty portal-field around Monolith (or a Star), Wave is oblationed.

\clearpage % ..........................................................

% \vspace*{1.1\baselineskip}
\subsubsection*{Teleporting Wave}
\addcontentsline{toc}{subsubsection}{Teleporting Wave}
\label{sec:One/Starchild/Activating Wave/Teleporting Wave}

\vspace*{-0.9\baselineskip}
\noindent
\begin{figure}[!h]
\includegraphics[width=1.0\textwidth, keepaspectratio=true]{examples/22_o/scn_o_19_star_moved_wave_teleportation.png}
\caption{Optional Wave teleportation}
\label{fig:scn_o_19_star_moved_wave_teleportation}
\end{figure}

\vspace*{-0.3\baselineskip}
Wave activated by pieces other than Starchild can still teleport as usual. Stars
in this variant can be moved out of their default positions. Teleportation for
Wave reaching a Star is optional, step-fields behind a Star are still accessible.
Here, light Wave could also activate light Queen. So,
\hyperref[fig:scn_d_10_teleport_wave_via_monolith]{Monolith is the only piece Wave cannot "pass-through"},
i.e. ignore as all the other pieces on chessboard.

\clearpage % ..........................................................

\vspace*{-2.1\baselineskip}
\noindent
\begin{figure}[!h]
\includegraphics[width=1.0\textwidth, keepaspectratio=true]{examples/22_o/scn_o_20_star_moved_wave_off_board.png}
\caption{Wave teleported off-board}
\label{fig:scn_o_20_star_moved_wave_off_board}
\end{figure}

Wave can end up with all step-fields off-board after teleportation, due to one or
both Stars moved out of their initial positions. In such a case, Wave is oblationed,
the same as in
\hyperref[fig:scn_d_12_wave_teleported_off_board]{previous variant, Discovery}.

Wave is also removed from chessboard if, after teleportation, all of its step-fields
are blocked; this is again similar to
\hyperref[fig:scn_d_11_teleported_wave_blocked]{previous variant, Discovery}.

\clearpage % ..........................................................

% \vspace*{1.1\baselineskip}
\subsubsection*{Steps after teleportation}
\addcontentsline{toc}{subsubsection}{Steps after teleportation}
\label{sec:One/Starchild/Activating Wave/Steps after teleportation}

\vspace*{-1.5\baselineskip}
\noindent
\begin{figure}[!h]
\includegraphics[width=1.0\textwidth, keepaspectratio=true]{examples/22_o/scn_o_21_steps_after_teleport.png}
\vspace*{-1.4\baselineskip}
\caption{Steps after teleportation}
\label{fig:scn_o_21_steps_after_teleport}
\end{figure}

\vspace*{-0.5\baselineskip}
Wave, activated by Unicorn (or Centaur) after teleportation has to follow two initially
chosen steps (long and short jump) according to a color of step-fields. Similar to
\hyperref[fig:scn_d_15_steps_after_teleport_init]{previous example with Monoliths},
two same-color Stars can be moved onto opposite-color fields. So, a two-step pattern
after teleportation remains the same, but order of steps is reversed. Here, after
teleporting into Star A, Wave emerges from Star B using the same two initial steps,
only in reversed order.

% ----------------------------------------------------- Activating Wave
\clearpage % ..........................................................
% Trance-journey ------------------------------------------------------

\subsection*{Trance-journey}
\addcontentsline{toc}{subsection}{Trance-journey}
\label{sec:One/Starchild/Trance-journey}

\vspace*{-1.4\baselineskip}
\noindent
\begin{figure}[!h]
\includegraphics[width=1.0\textwidth, keepaspectratio=true]{examples/22_o/scn_o_22_light_starchild_init_trance_journey.png}
\caption{Light Starchild initiating}
\label{fig:scn_o_22_light_starchild_init_trance_journey}
\end{figure}

Trance-journey can be initiated by Starchild activating another Starchild;
colors of Starchilds do not need to match. If the same, activation can be
done on step-fields (here), otherwise only on miracle-fields (next example).
Initiating Starchild must activate entrancing Starchild directly, there must
be no pieces between the two.

\clearpage % ..........................................................

\vspace*{-2.1\baselineskip}
\noindent
\begin{figure}[!h]
\includegraphics[width=1.0\textwidth, keepaspectratio=true]{examples/22_o/scn_o_23_dark_starchild_init_trance_journey.png}
\caption{Dark Starchild initiating}
\label{fig:scn_o_23_dark_starchild_init_trance_journey}
\end{figure}

Entrancing Starchild can activate any piece in the same color (except Kings, Waves,
Starchilds, Stars and Monoliths), on its miracle-fields. Entranced piece can
then optionally take on trance-journey, or take an ordinary step with one momentum.
If entrancing Starchild is light, piece can take trance-journey in one chosen
direction from
\hyperref[fig:scn_cot_17_light_shaman_trance_journey]{light Shaman's pattern},
otherwise from
\hyperref[fig:scn_cot_19_dark_shaman_trance_journey]{dark Shaman's pattern}.

\clearpage % ..........................................................

\vspace*{-2.1\baselineskip}
\noindent
\begin{figure}[!h]
\includegraphics[width=1.0\textwidth, keepaspectratio=true]{examples/22_o/scn_o_24_shaman_init_trance_journey.png}
% \vspace*{-0.6\baselineskip} # does the opposite of instruction ???
\caption{Shaman initiating}
\label{fig:scn_o_24_shaman_init_trance_journey}
\end{figure}

\vspace*{-0.6\baselineskip}
Alternatively, a Shaman can activate entrancing Starchild, not necessarily in the same
colors. If initiated by Shaman, only Waves are permitted between it and Starchild.
Initiating piece, regardless if Starchild or Shaman, can be activated by other pieces,
including opponent's.

Entranced piece can take on trance-journey even if it received no momentum;
length of trance-journey does not depend on received momentum.

\clearpage % ..........................................................

\vspace*{-2.1\baselineskip}
\noindent
\begin{figure}[!h]
\includegraphics[width=1.0\textwidth, keepaspectratio=true]{examples/22_o/scn_o_25_trance_journey_started_by_shaman.png}
\caption{Light-pattern trance-journey}
\label{fig:scn_o_25_trance_journey_started_by_shaman}
\end{figure}

Entranced piece can end its trance-journey on any empty step-field, of chosen
direction (color). If all are occupied, then it emerges on any empty entrancing
Starchild's miracle-field. If there's none, then it emerges on empty initiating
Shaman's or Starchild's miracle-fields. If all are occupied, entranced piece is
oblationed.

\clearpage % ..........................................................

% \vspace*{-1.1\baselineskip}
\subsubsection*{Push-pull entrancement}
\addcontentsline{toc}{subsubsection}{Push-pull entrancement}
\label{sec:One/Starchild/Trance-journey/Push-pull entrancement}

\noindent
\begin{minipage}{\textwidth}
\begin{wrapfigure}[7]{l}{0.4\textwidth}
\centering
\includegraphics[width=0.346153846\textwidth, keepaspectratio=true]{examples/22_o/scn_o_26_push_pull_trance_journey_init.png}
\caption{Initiating trance-journey}
\label{fig:scn_o_26_push_pull_trance_journey_init}
\end{wrapfigure}
Starchild initiating trance-journey could also be activated later in the same cascade,
and act as an entrancing Shaman. This is similar to push-pull entrancement in the
\hyperref[fig:star/scn_cot_36_push_pull_entrancement_start]{previous variant, Conquest of Tlalocan}.
\end{minipage}

\vspace*{4.9\baselineskip}
\noindent
\begin{minipage}{\textwidth}
\begin{wrapfigure}[10]{l}{0.4\textwidth}
\centering
\includegraphics[width=0.346153846\textwidth, keepaspectratio=true]{examples/22_o/scn_o_27_push_pull_trance_journey_entrancing.png}
\caption{Push-pull entrancing}
\label{fig:scn_o_27_push_pull_trance_journey_entrancing}
\end{wrapfigure}
In previous eample, dark Starchild activated Wave A, which then activated Wave B.
Here, Wave B is "in the air", about to activate dark Starchild, which will then
entrance dark Knight.
\end{minipage}

\clearpage % ..........................................................

\vspace*{-2.1\baselineskip}
\noindent
\begin{figure}[!h]
\includegraphics[width=1.0\textwidth, keepaspectratio=true]{examples/22_o/scn_o_28_push_pull_trance_journey_entranced.png}
\caption{Dark-pattern trance-journey}
\label{fig:scn_o_28_push_pull_trance_journey_entranced}
\end{figure}

Starchild can initiate trance-journey by
\hyperref[sec:Terms/Push-pull activation]{push-pull activation},
if its color is the same as color of entranced piece; here both Starchild and Knight
are dark.

Push-pull activation would work even if initiating Starchild has been activated by
some other pieces, which don't have to be in the same color.

\clearpage % ..........................................................

% \vspace*{-1.1\baselineskip}
\subsubsection*{Failed trance-journey}
\addcontentsline{toc}{subsubsection}{Failed trance-journey}
\label{sec:One/Starchild/Trance-journey/Failed trance-journey}

\vspace*{-1.1\baselineskip}
\noindent
\begin{figure}[!h]
\includegraphics[width=1.0\textwidth, keepaspectratio=true]{examples/22_o/scn_o_29_trance_journey_failed.png}
\caption{Failed trance-journey}
\label{fig:scn_o_29_trance_journey_failed}
\end{figure}

If all step-fields in a
\hyperref[fig:scn_cot_12_trance_journey_init]{regular trance-journey}
are blocked by Kings, Stars or Monoliths, entranced Shaman is
\hyperref[sec:Terms/Oblation]{oblationed}, i.e. removed from chessboard
as if captured by opponent.

For a trance-journey to fail, since it's optional, all capture-fields needs
to be either empty or blocked by own pieces, and all step-fields of a normal
movement also needs to be blocked.

\clearpage % ..........................................................

\vspace*{-2.1\baselineskip}
\noindent
\begin{figure}[!h]
\includegraphics[width=1.0\textwidth, keepaspectratio=true]{examples/22_o/scn_o_30_trance_journey_failed_2.png}
\caption{Failed new trance-journey}
\label{fig:scn_o_30_trance_journey_failed_2}
\end{figure}

In addition to all step-fields being blocked, for a Starchild-induced trance-journey
to fail, all miracle-fields of initiating piece and entrancing Starchild needs
to be blocked as well.

In contrast to
\hyperref[fig:scn_o_29_trance_journey_failed]{failed Shaman's trance-journey},
Starchild-induced trance-journey is blocked by any piece on its step-fields;
here, own and opponent's Pawns.

% ------------------------------------------------------ Trance-journey
\clearpage % ..........................................................
% Syzygy --------------------------------------------------------------

\subsection*{Syzygy}
\addcontentsline{toc}{subsection}{Syzygy}
\label{sec:One/Starchild/Syzygy}

\vspace*{-1.4\baselineskip}
\noindent
\begin{figure}[!h]
\includegraphics[width=1.0\textwidth, keepaspectratio=true]{examples/22_o/scn_o_31_syzygy_monolith.png}
% \vspace*{-1.5\baselineskip}
\caption{Demoting-to-Pawn syzygy}
\label{fig:scn_o_31_syzygy_monolith}
\end{figure}

% \vspace*{-0.7\baselineskip}
Starchild is celestial piece, it can participate in
\hyperref[fig:scn_d_19_syzygy_2_stars_init]{demoting-to-Pawn syzygy} in place of
Stars and Monoliths. Starchilds are also eligible to demotion.

\hyperref[fig:scn_d_18_syzygy_explain]{Again}, shortest step connecting
Stars, Monoliths, Starchilds is called syzygy-step, fields which are connected by
syzygy-steps are called syzygy-fields.

\clearpage % ..........................................................
% Resurrection syzygy .................................................

\subsubsection*{Resurrection syzygy}
\addcontentsline{toc}{subsubsection}{Resurrection syzygy}
\label{sec:One/Starchild/Syzygy/Resurrection syzygy}

\vspace*{-1.4\baselineskip}
\noindent
\begin{figure}[!h]
\includegraphics[width=1.0\textwidth, keepaspectratio=true]{examples/22_o/scn_o_32_syzygy_starchild_init.png}
\caption{Resurrection syzygy start}
\label{fig:scn_o_32_syzygy_starchild_init}
\end{figure}

If Starchild moves into syzygy, from a field not in any alignment, onto an empty
syzygy-field, it's granted option to resurrect one captured piece. A piece is
resurrected by replacing initiating Starchild, Starchild itself is then oblationed.
Only captured pieces can be resurrected. Kings, Stars and Monoliths cannot be
resurrected.

\clearpage % ..........................................................

\vspace*{-2.1\baselineskip}
\noindent
\begin{figure}[!h]
\includegraphics[width=1.0\textwidth, keepaspectratio=true]{examples/22_o/scn_o_33_syzygy_starchild_end.png}
\vspace*{-1.4\baselineskip}
\caption{Queen resurrected}
\label{fig:scn_o_33_syzygy_starchild_end}
\end{figure}

\vspace*{-0.3\baselineskip}
Here, resurrected Queen replaced initiating Starchild. Note, in this variant
\hyperref[sec:One/Promotion]{promotion is monogamous}, so the only light Queen
had to be captured, before it could be resurrected.

Resurrecting is not limited to player's own pieces, opponent's pieces can be
resurrected as well. For instance, dark Queen could be resurrected instead of
the light one. Again, only captured pieces can be resurrected. In this variant,
to resurrect Queen, a Pawn must not be promoted to one.

\clearpage % ..........................................................

\vspace*{-2.1\baselineskip}
\noindent
\begin{figure}[!h]
\includegraphics[width=1.0\textwidth, keepaspectratio=true]{examples/22_o/scn_o_34_syzygy_starchild_resurrection.png}
\caption{Starchild resurrected}
\label{fig:scn_o_34_syzygy_starchild_resurrection}
\end{figure}

Captured Waves and Starchilds can be resurrected, without initiating Starchild
being oblationed. Chosen piece can emerge on any empty miracle-field around
Starchilds in syzygy (here, fields marked 1). If miracle-fields are all
occupied, piece emerges on any empty portal-field around Stars, Monoliths in
syzygy (fields 2). If all portal-fields are occupied, piece emerges on any
empty syzygy-field (3). If all are occupied, resurrection is not performed.

% ................................................. Resurrection syzygy
\clearpage % ..........................................................

\subsubsection*{Reentering syzygy}
\addcontentsline{toc}{subsubsection}{Reentering syzygy}
\label{sec:One/Starchild/Syzygy/Reentering syzygy}

\vspace*{-1.2\baselineskip}
\noindent
\begin{figure}[!h]
\includegraphics[width=1.0\textwidth, keepaspectratio=true]{examples/22_o/scn_o_35_reentering_syzygies.png}
\caption{Reentering syzygy}
\label{fig:scn_o_35_reentering_syzygies}
\end{figure}

Similar to
\hyperref[fig:scn_d_23_syzygy_reentering_same_move]{Monolith reentering syzygy},
to get option to resurrect captured piece, Starchild has to move from a normal,
non-syzygy field into syzygy. Starchild already in syzygy can move into the same
(or the other) alignment, but cannot resurrect any piece. To be able to resurrect,
Starchild has to move out of alignment in a first move, and then on a next move
it can move into resurrecting syzygy.

\clearpage % ..........................................................

\subsubsection*{Cascading syzygy}
\addcontentsline{toc}{subsubsection}{Cascading syzygy}
\label{sec:One/Starchild/Syzygy/Cascading syzygy}

\vspace*{-1.2\baselineskip}
\noindent
\begin{figure}[!h]
\includegraphics[width=1.0\textwidth, keepaspectratio=true]{examples/22_o/scn_o_36_syzygy_starchild_cascading.png}
\caption{Starchild cascading}
\label{fig:scn_o_36_syzygy_starchild_cascading}
\end{figure}

Starchild moving into syzygy can activate pieces on a destination syzygy-field,
and start a cascade; this move does not grant option to resurrect. To be able to
resurrect, Starchild has to move onto empty syzygy-field, from a field out of
any alignment.

\clearpage % ..........................................................

\subsubsection*{Double syzygy}
\addcontentsline{toc}{subsubsection}{Double syzygy}
\label{sec:One/Starchild/Syzygy/Double syzygy}

\vspace*{-1.2\baselineskip}
\noindent
\begin{figure}[!h]
\includegraphics[width=1.0\textwidth, keepaspectratio=true]{examples/22_o/scn_o_38_two_syzygies_shared_celestial_piece.png}
\caption{Shared celestial piece}
\label{fig:scn_o_38_two_syzygies_shared_celestial_piece}
\end{figure}

Starchild initiating syzygy can activate another Starchild, creating double
(triple, ...) syzygies in the same move; only the very last Starchild in a
cascade can resurrect a captured piece, if it stepped onto an empty syzygy-field.
This is so regardless if syzygies share celestial piece (pictured here), just a
syzygy-field, or are completely independent.

\clearpage % ..........................................................

\subsubsection*{Opponent's Starchild}
\addcontentsline{toc}{subsubsection}{Opponent's Starchild}
\label{sec:One/Starchild/Syzygy/Opponent's Starchild}

\vspace*{-1.2\baselineskip}
\noindent
\begin{figure}[!h]
\includegraphics[width=1.0\textwidth, keepaspectratio=true]{examples/22_o/scn_o_41_syzygy_opponents_starchild.png}
\caption{Opponent's Starchild in syzygy}
\label{fig:scn_o_41_syzygy_opponents_starchild}
\end{figure}

Opponent's Starchild can be activated, and pushed into syzygy, this will also
initiate resurrection. Here, light player is moving dark Starchild into syzygy.

Any Starchild can resurrect any captured piece, regardless of their respective
colors. Here, light player can resurrect any captured piece, regardless if it's
light or dark.

\clearpage % ..........................................................

\subsubsection*{Star-initiated syzygy}
\addcontentsline{toc}{subsubsection}{Star-initiated syzygy}
\label{sec:One/Starchild/Syzygy/Star-initiated syzygy}

\vspace*{-1.2\baselineskip}
\noindent
\begin{figure}[!h]
\includegraphics[width=1.0\textwidth, keepaspectratio=true]{examples/22_o/scn_o_42_syzygy_star.png}
\caption{Star-initiated syzygy}
\label{fig:scn_o_42_syzygy_star}
\end{figure}

Pushing Star into a syzygy brings no additional interactions, i.e. it neither can
demote own figure to Pawn, nor it can resurrect any captured piece.

% -------------------------------------------------------------- Syzygy
\clearpage % ..........................................................
% Divergence ----------------------------------------------------------

\subsection*{Divergence}
\addcontentsline{toc}{subsection}{Divergence}
\label{sec:One/Starchild/Divergence}

\vspace*{-1.2\baselineskip}
\noindent
\begin{figure}[!h]
\includegraphics[width=1.0\textwidth, keepaspectratio=true]{examples/22_o/scn_o_43_starchild_divergence_init.png}
\caption{Own Starchild is divergent}
\label{fig:scn_o_43_starchild_divergence_init}
\end{figure}

. . .

% \TODO

% \clearpage % ..........................................................

\vspace*{-2.1\baselineskip}
\noindent
\begin{figure}[!h]
\includegraphics[width=1.0\textwidth, keepaspectratio=true]{examples/22_o/scn_o_44_starchild_divergence_end.png}
\caption{Diverging Queen}
\label{fig:scn_o_44_starchild_divergence_end}
\end{figure}

. . .

\TODO

% \clearpage % ..........................................................

% % \vspace*{1.1\baselineskip}
% \subsubsection*{Starchild is divergent}
% \addcontentsline{toc}{subsubsection}{Starchild is divergent}
% \label{sec:One/Starchild/Movement/Starchild is divergent}

% \clearpage % ..........................................................

% % \vspace*{1.1\baselineskip}
% \subsubsection*{Starchild is not transparent}
% \addcontentsline{toc}{subsubsection}{Starchild is not transparent}
% \label{sec:One/Starchild/Movement/Starchild is not transparent}

% \TODO

\clearpage % ..........................................................

% \vspace*{1.1\baselineskip}
\subsubsection*{Starchild cannot diverge}
\addcontentsline{toc}{subsubsection}{Starchild cannot diverge}
\label{sec:One/Starchild/Movement/Starchild cannot diverge}

\TODO

% ---------------------------------------------------------- Divergence
% *********************************************************** Starchild
\clearpage % ..........................................................

\section*{Rush, en passant}
\addcontentsline{toc}{section}{Rush, en passant}
\label{sec:One/Rush, en passant}

\vspace*{-1.2\baselineskip}
\noindent
\begin{figure}[!h]
\includegraphics[width=1.0\textwidth, keepaspectratio=true]{en_passants/22_one_en_passant.png}
\caption{En passant}
\label{fig:22_one_en_passant}
\end{figure}

Rush and en passant are identical to those in
\hyperref[fig:14_hemera_s_dawn_en_passant]{Hemera's Dawn variant}.
Own Pawns can be rushed for up to 11 fields in this variant.

\clearpage % ..........................................................

\section*{Promotion}
\addcontentsline{toc}{section}{Promotion}
\label{sec:One/Promotion}

Promotion is non enforced, delayed variety, i.e. it's the same as in
\hyperref[sec:Age of Aquarius/Promotion]{previous chess variant}, Age of Aquarius.

Additionaly, promotion in this variant is monogamous. Only one Queen
in the same color can be present on chessboard at any given time.

\hyperref[fig:scn_n_22_only_one_queen]{Similarly to previous variant, Nineteen},
if own Queen is present on a chessboard, opponent's Queen cannot be
converted, and has to be captured instead.

\clearpage % ..........................................................

\section*{Castling}
\addcontentsline{toc}{section}{Castling}
\label{sec:One/Castling}

Castling is
\hyperref[sec:Nineteen/Castling]{the same as in Nineteen variant},
only difference is that King can move
between 2 and 10 fields across. All other constraints from Nineteen variant still
applies.

\noindent
\begin{figure}[!h]
\includegraphics[width=1.0\textwidth, keepaspectratio=true]{castlings/22_o/one_castling.png}
\caption{Castling}
\label{fig:one_castling}
\end{figure}

In example above, all valid King's castling moves are numbered.

\noindent
\begin{figure}[!h]
\includegraphics[width=1.0\textwidth, keepaspectratio=true]{castlings/22_o/one_castling_right_04.png}
\caption{Castling short right}
\label{fig:one_castling_right_04}
\end{figure}

In this example King was castling short to the right. Initial King's position is
marked with "K". After castling is finished, right Rook ends up at field immediately
left to the King.

\clearpage % ..........................................................

\section*{Initial setup}
\addcontentsline{toc}{section}{Initial setup}
\label{sec:One/Initial setup}

Compared to initial setup of Discovery, Starchild is inserted between Centaur and
Serpent symmetrically, on both sides of chessboard. This can be seen in the image
below:

\noindent
\begin{figure}[h]
\includegraphics[width=1.0\textwidth, keepaspectratio=true]{boards/22_one.png}
\caption{One board}
\label{fig:22_one}
\end{figure}

\clearpage % ..........................................................
% ========================================================= One chapter
